\section{Runtime Evolution Around a Fixed Model}
\label{sec:evolution}\label{sec:runtimeevolution}

\begin{designbox}
Let each mission update the operating state \emph{around} a fixed model --- memory,
skills, tools, verifiers, routing, and evaluations --- rather than the model's
weights. Make every update attributable: each change names its \emph{source} and
the \emph{authoritative owner} that commits it to a named persistence surface. That
owner is often a role distinct from the one that did the work, but not always ---
some components are owner-complete by design. Not every component moves every run,
and none of it is a gradient step on the base model.
\end{designbox}

Act~II showed the machine that turns one mission into an audited increment. Act~III
states what that increment \emph{is} and how it accumulates. The report's master
spine (Figure~\ref{fig:spine}) names the step as $H(t{+}1)=U(H(t),\text{trajectory},
\text{evidence})$; here we make $H$ concrete from the checkpoint, skills, wiki,
tools, verifiers, routing, and evaluation facts already established
(Sections~\ref{sec:rolesplanes}, \ref{sec:lifecycle}, \ref{sec:evidence}). Each
mission transforms that state under the evidence it accepted.

\begin{equation}
H_t = \{M_t,S_t,A_t,V_t,R_t,Q_t\},
\qquad
H_{t+1}=U(H_t,\tau_t,E_t),
\qquad
\theta_{t+1}=\theta_t .
\end{equation}

Here $M$ is memory --- the curated \code{checkpoint.json} and the recency journal
derived from the event tape; $S$ is the versioned skill library; $A$ is the
available tool and vertical-procedure surface; $V$ is the verifier and
evidence-procedure set --- the per-stage checklist and the anti-cheat construction;
$R$ is routing and role configuration --- backend/model selection and campaign
lifetime; and $Q$ is the evaluation and task-definition set --- the backlog and its
scopes. The mission contributes a process trajectory $\tau_t$ and the accepted
evidence $E_t$, and $U$ folds them into the next state. The final equality carries
the design commitment: $\theta_{t+1}=\theta_t$. The weights of the underlying
coding-and-reasoning model are \emph{fixed} for the campaign. Runtime evolution is
\textbf{not online parameter training} --- no gradient step is taken on the base
model, and no claim of weight-level learning is made anywhere in this report. What
changes is the operating state the fixed model reads and writes.

\subsection{Every update is attributable}

Runtime evolution is not one mechanism but a family, already introduced as four
ownership paths --- automatic, model-proposed, Reviewer-authored, and
operator-invoked (Section~\ref{sec:life}). Table~\ref{tab:evolution} resolves that
per component: the \emph{source} of a change to $H$, the \emph{authoritative owner}
that commits it, and the named surface where it persists. For the two components the
system \emph{earns} during a mission --- memory ($M$) and skills ($S$) --- the owner
is deliberately a role other than the one that produced the work: the Reviewer
certifies a checkpoint and skill edits it did not author itself, the same
work-versus-certification separation that governs completion. The remaining
components are owner-complete by design and honestly labeled as such. Verifiers
($V$) --- the per-stage checklist --- are \textbf{Planner-owned}: the Planner both
authors and applies \code{checklist\_ops} (\srcpath{planner/planner.py}), and the
Reviewer is \textbf{feedback-only}, reporting a wrong checklist through its
\code{checklist\_feedback} channel (\srcpath{reviewer/\_core.py}) rather than editing
it; the external check on a Planner-authored checklist is that feedback plus the
Reviewer's independent completion authority, not an accept/commit gate. Tools and
procedures ($A$) are \textbf{operator-owned} --- a human outside the four-role loop
installs and retires them. Routing ($R$) is operator-sourced and Manager-committed;
evaluations ($Q$) are Planner-authored and scheduler-committed.

\begin{table}[t]
\centering
\small
\begin{tabularx}{\linewidth}{@{}p{3.0cm} p{2.3cm} p{2.45cm} X@{}}
\toprule
\textbf{State component} & \textbf{Change source} & \textbf{Authoritative owner} & \textbf{Persistence surface} \\
\midrule
$M$ --- memory & Engineer & Reviewer & \code{checkpoint.json}; event-tape journal \\
$S$ --- skills & Engineer / Scientist & Reviewer & \code{skills/} versioned Markdown (\code{skill\_ops}) \\
$A$ --- tools / procedures & operator / vertical & operator & installed tool surface; vertical stages \\
$V$ --- verifiers & Planner & Planner & per-stage checklist (\code{checklist\_ops}); Reviewer \code{checklist\_feedback} \\
$R$ --- routing / roles & operator & Manager & \code{continuous.json}; knob store (\code{CONFIG}) \\
$Q$ --- evaluations / tasks & Planner & scheduler & \code{backlog.jsonl} (\code{TaskSpec}) \\
\bottomrule
\end{tabularx}
\caption{The runtime-state components of $H$: the source of each change, the
authoritative owner that commits it, and the one persistence surface where it lands
(with the interface object that carries it). For $M$ and $S$ the owner is a role
distinct from the one that did the work, so the Reviewer certifies output it did not
produce; $V$ is Planner-owned (the Planner authors and applies \code{checklist\_ops};
the Reviewer's \code{checklist\_feedback} is an external report, not a commit gate),
and $A$ is operator-owned. Sources: Sections~\ref{sec:rolesplanes} and
\ref{sec:lifecycle}.}
\label{tab:evolution}
\end{table}

Two disciplines keep the equation honest. First, $U$ is \emph{partial}: a typical
mission touches only a subset of $H$. Many missions add no new skill, synthesize no
wiki page, and leave tools, routing, and evaluations untouched --- $A$ in particular
is the most static component, changing mainly by operator or code action. The state
is updated where accepted evidence warrants it, not rewritten wholesale each run.
Second, because $\theta$ is fixed, all durable capability lives in $H$ on disk, where
it is inspectable and reversible --- not latent in model weights. That is what lets
the next section treat the \emph{process}, not only the final artifact, as the asset
the system actually retains.
